% ================================================================
% MBP-ALC — Preámbulo LaTeX (XeLaTeX)
% Estilo moderno y académico para libro técnico CEPAL
% ================================================================

% --- Fuentes via fontspec (XeLaTeX) — fonts del sistema Windows --
% Cambria / Calibri / Consolas son fonts nativas de Windows;
% no requieren instalación adicional en TinyTeX.
% Para usar TeX Gyre (calidad superior), instalar primero:
%   tinytex::tlmgr_install(c("tex-gyre", "tex-gyre-math"))
% y reemplazar los tres \set*font por los bloques comentados abajo.
\usepackage{fontspec}

\setmainfont{Cambria}[
  BoldFont       = {Cambria Bold},
  ItalicFont     = {Cambria Italic},
  BoldItalicFont = {Cambria Bold Italic}
]
% Alternativa de mayor calidad (requiere instalación):
% \setmainfont{TeX Gyre Pagella}[...]

\setsansfont{Calibri}[
  Scale          = MatchLowercase,
  BoldFont       = {Calibri Bold},
  ItalicFont     = {Calibri Italic},
  BoldItalicFont = {Calibri Bold Italic}
]
% \setsansfont{TeX Gyre Heros}[...]

\setmonofont{Consolas}[Scale=0.88]
% \setmonofont{Inconsolata}[Scale=0.88]

% --- Matemáticas ------------------------------------------------
% Cambria Math es el motor matemático nativo de Windows y está
% siempre disponible; empareja visualmente con Cambria.
\usepackage{unicode-math}
\setmathfont{Cambria Math}
% \setmathfont{TeX Gyre Pagella Math}   % alternativa tras instalar tex-gyre-math
\usepackage{amsmath}
\usepackage{amssymb}
\usepackage{amsthm}
\usepackage{mathtools}

% --- Idioma -----------------------------------------------------
% polyglossia/babel son cargados automáticamente por bookdown
% a partir de lang: es en el YAML de index.Rmd.
% No redefinir aquí para evitar el conflicto "\textspanish already defined".

% --- Colores del libro ------------------------------------------
\usepackage{xcolor}
\definecolor{cepalnavy}{RGB}{27,  42,  74}   % azul marino CEPAL
\definecolor{cepalblue}{RGB}{37,  99, 235}   % azul acento
\definecolor{cepalgray}{RGB}{74,  85, 104}   % gris texto secundario
\definecolor{cepalltbg}{RGB}{248,250,252}    % fondo código/tabla
\definecolor{cepalborder}{RGB}{226,232,240}  % borde suave

% --- Geometría de página ----------------------------------------
\usepackage[
  top      = 2.5cm,
  bottom   = 3.0cm,
  inner    = 3.0cm,
  outer    = 2.2cm,
  headheight = 14pt,
  headsep    = 0.6cm
]{geometry}

% --- Interlineado -----------------------------------------------
\usepackage{setspace}
\setstretch{1.15}

% --- Estilo de capítulos y secciones ----------------------------
\usepackage{titlesec}

% Capítulo: número grande a la derecha, título con línea inferior
\titleformat{\chapter}[display]
  {\normalfont\sffamily\huge\bfseries\color{cepalnavy}}
  {%
    \raggedleft
    \textcolor{cepalborder}{\fontsize{72}{72}\selectfont\thechapter}%
  }
  {-1.2ex}
  {%
    \vspace{0.8ex}%
    \textcolor{cepalblue}{\rule{\textwidth}{1.5pt}}%
    \vspace{0.6ex}%
  }
  [\vspace{0.3ex}\textcolor{cepalborder}{\rule{\textwidth}{0.4pt}}]

\titleformat{\section}
  {\normalfont\sffamily\Large\bfseries\color{cepalnavy}}
  {\thesection}{1em}{}

\titleformat{\subsection}
  {\normalfont\sffamily\large\bfseries\color{cepalgray}}
  {\thesubsection}{1em}{}

\titleformat{\subsubsection}
  {\normalfont\sffamily\normalsize\bfseries\color{cepalgray}}
  {\thesubsubsection}{1em}{}

% Espaciado de secciones
\titlespacing*{\chapter}    {0pt}{0pt}{2.8em}
\titlespacing*{\section}    {0pt}{2.0em}{0.6em}
\titlespacing*{\subsection} {0pt}{1.5em}{0.4em}

% --- Encabezados y pies de página -------------------------------
\usepackage{fancyhdr}
\pagestyle{fancy}
\fancyhf{}
\renewcommand{\headrulewidth}{0.4pt}
\renewcommand{\footrulewidth}{0pt}
\renewcommand{\headrule}{\color{cepalborder}\hrule width\headwidth height\headrulewidth}

\fancyhead[LE]{\small\sffamily\color{cepalgray}\leftmark}
\fancyhead[RO]{\small\sffamily\color{cepalgray}\rightmark}
\fancyfoot[C]{\small\sffamily\color{cepalgray}\thepage}

\fancypagestyle{plain}{
  \fancyhf{}
  \renewcommand{\headrulewidth}{0pt}
  \fancyfoot[C]{\small\sffamily\color{cepalgray}\thepage}
}

% --- Tablas -----------------------------------------------------
\usepackage{booktabs}
\usepackage{longtable}
\usepackage{array}
\usepackage{colortbl}
\usepackage{multirow}
\usepackage{makecell}

% Encabezado de tabla con color CEPAL
\newcommand{\tablehead}[1]{%
  \rowcolor{cepalnavy}\color{white}\bfseries\sffamily #1%
}

% --- Código / bloques de sintaxis -------------------------------
\usepackage{fancyvrb}
\usepackage{fvextra}

\DefineVerbatimEnvironment{Highlighting}{Verbatim}{
  commandchars = \\\{\},
  fontsize     = \small,
  breaklines,
  breakanywhere,
  frame        = leftline,
  framerule    = 3pt,
  rulecolor    = \color{cepalblue},
  framesep     = 10pt,
  fillcolor    = \color{cepalltbg},
  numbers      = none
}

% --- Cajas especiales (notas, definiciones, etc.) ---------------
\usepackage[most, breakable]{tcolorbox}

\tcbset{
  base style/.style={
    enhanced, breakable,
    fonttitle = \sffamily\bfseries\small,
    coltitle  = white,
    arc       = 4pt,
    boxrule   = 0.4pt,
    left      = 8pt,
    right     = 8pt,
    top       = 6pt,
    bottom    = 6pt,
  }
}

\newtcolorbox{nota}[1][Nota]{
  base style,
  title       = #1,
  colback     = cepalltbg,
  colframe    = cepalblue,
  colbacktitle= cepalblue
}

\newtcolorbox{definicion}[1][Definición]{
  base style,
  title        = #1,
  colback      = cepalltbg,
  colframe     = cepalnavy,
  colbacktitle = cepalnavy
}

\newtcolorbox{advertencia}[1][Advertencia]{
  base style,
  title        = #1,
  colback      = {yellow!8!white},
  colframe     = {orange!70!black},
  colbacktitle = {orange!80!black}
}

% --- Figuras y leyendas -----------------------------------------
\usepackage{caption}
\captionsetup{
  font        = {small, sf},
  labelfont   = {bf, color=cepalnavy},
  labelsep    = period,
  justification = centering,
  skip        = 6pt
}

% --- Hiperenlaces -----------------------------------------------
\usepackage{hyperref}
\hypersetup{
  colorlinks  = true,
  linkcolor   = cepalnavy,
  citecolor   = cepalblue,
  urlcolor    = cepalblue,
  pdftitle    = {Modelos Bayesianos de Población para ALC},
  pdfauthor   = {Andrés Gutiérrez y Stalyn Guerrero},
  pdfsubject  = {Estadística demográfica bayesiana},
  pdfkeywords = {modelos bayesianos, población, América Latina, estimación subnacional}
}

% --- Listas -----------------------------------------------------
\usepackage{enumitem}
\setlist{itemsep = 0.25em, parsep = 0.15em, topsep = 0.4em}
\setlist[itemize]{leftmargin = 1.5em}
\setlist[enumerate]{leftmargin = 1.8em}

% --- Algoritmos -------------------------------------------------
\usepackage[ruled, vlined, linesnumbered, spanish]{algorithm2e}

% --- Índice y listas de figuras/tablas --------------------------
\usepackage{tocloft}
\renewcommand{\cftchapfont}{\sffamily\bfseries\color{cepalnavy}}
\renewcommand{\cftsecfont}{\sffamily}
\renewcommand{\cftsubsecfont}{\sffamily\small}
\renewcommand{\cftchappagefont}{\sffamily\bfseries\color{cepalnavy}}
\setlength{\cftbeforechapskip}{0.6em}
